\documentclass[12pt]{article}
\usepackage{amsmath}
\usepackage{enumitem}
\usepackage[margin=1in]{geometry}

\title{DNS and HTTP Practice Problems}
\author{}
\date{}

\begin{document}

\maketitle

\section*{DNS Practice Problems}

\subsection*{Problem 1: DNS Resolution Chain}
A user types:
\[
\texttt{www.eng.bu.edu}
\]
into a browser. Explain the DNS resolution process.

\subsection*{Solution}
The user's laptop first sends the query to a recursive resolver, such as an ISP resolver, Google DNS, or Cloudflare DNS.

If the recursive resolver does not already have the answer cached, it performs iterative queries:

\begin{enumerate}
    \item Ask a root DNS server:
    \[
    \text{Who handles } \texttt{.edu}?
    \]
    \item Ask a \texttt{.edu} TLD server:
    \[
    \text{Who handles } \texttt{bu.edu}?
    \]
    \item Ask the \texttt{bu.edu} authoritative server:
    \[
    \text{Who handles } \texttt{eng.bu.edu}?
    \]
    \item Ask the \texttt{eng.bu.edu} authoritative server:
    \[
    \text{What is the A record for } \texttt{www.eng.bu.edu}?
    \]
\end{enumerate}

The final IP address is returned to the recursive resolver, then back to the user's laptop.

\[
\boxed{\text{The laptop only talks directly to the recursive resolver.}}
\]

\subsection*{Problem 2: DNS Record Types}
Match each DNS record type with its purpose.

\[
\begin{array}{c|c}
\text{Record Type} & \text{Purpose} \\
\hline
A & ? \\
AAAA & ? \\
CNAME & ? \\
MX & ? \\
NS & ? \\
TXT & ?
\end{array}
\]

\subsection*{Solution}

\[
\begin{array}{c|c}
\text{Record Type} & \text{Purpose} \\
\hline
A & \text{Maps a domain name to an IPv4 address} \\
AAAA & \text{Maps a domain name to an IPv6 address} \\
CNAME & \text{Creates an alias to another domain name} \\
MX & \text{Specifies mail servers for a domain} \\
NS & \text{Specifies name servers for a zone} \\
TXT & \text{Stores text data such as SPF, DKIM, or verification records}
\end{array}
\]

\[
\boxed{\text{DNS records tell resolvers what kind of information a domain has.}}
\]

\subsection*{Problem 3: DNS Caching and TTL}
A DNS record has a TTL of:
\[
300 \text{ seconds}
\]
A recursive resolver receives the record at time:
\[
t = 0
\]
At what time must the resolver stop using the cached answer?

\subsection*{Solution}
TTL stands for Time To Live. It tells the resolver how long it may cache the answer.

\[
t_{\text{expire}} = t_{\text{start}} + TTL
\]

\[
t_{\text{expire}} = 0 + 300
\]

\[
\boxed{t_{\text{expire}} = 300 \text{ seconds}}
\]

After 300 seconds, the resolver must query again instead of using the cached record.

\subsection*{Problem 4: Low TTL vs High TTL}
A website sets its DNS TTL to 300 seconds instead of 3600 seconds. What is the tradeoff?

\subsection*{Solution}
A lower TTL means DNS changes update faster.

\[
300 \text{ seconds} = 5 \text{ minutes}
\]

\[
3600 \text{ seconds} = 1 \text{ hour}
\]

With a TTL of 300 seconds, users will receive updated DNS information sooner.

However, a lower TTL also means recursive resolvers must query authoritative servers more often.

\[
\boxed{\text{Low TTL gives faster updates but creates more DNS query load.}}
\]

\section*{HTTP Practice Problems}

\subsection*{Problem 5: HTTP Request}
Write a basic HTTP/1.1 GET request for:
\[
\texttt{https://example.com/index.html}
\]

\subsection*{Solution}
A basic HTTP/1.1 request looks like:

\[
\begin{array}{l}
\texttt{GET /index.html HTTP/1.1} \\
\texttt{Host: example.com} \\
\texttt{User-Agent: curl/8.4.0} \\
\texttt{Accept: */*}
\end{array}
\]

\[
\boxed{\text{The request line includes the method, path, and HTTP version.}}
\]

\subsection*{Problem 6: HTTP Response}
Given this HTTP response:

\[
\begin{array}{l}
\texttt{HTTP/1.1 200 OK} \\
\texttt{Content-Type: text/html; charset=UTF-8} \\
\texttt{Content-Length: 1256} \\
\texttt{Connection: keep-alive}
\end{array}
\]

Answer the following:
\begin{enumerate}
    \item Was the request successful?
    \item What type of content was returned?
    \item How large is the body?
    \item Will the TCP connection stay open?
\end{enumerate}

\subsection*{Solution}

\begin{enumerate}
    \item Yes. The status code is:
    \[
    200
    \]
    which means success.

    \item The content type is:
    \[
    \texttt{text/html}
    \]
    so the response contains an HTML document.

    \item The body length is:
    \[
    1256 \text{ bytes}
    \]

    \item Yes. The header:
    \[
    \texttt{Connection: keep-alive}
    \]
    means the TCP connection can be reused.
\end{enumerate}

\[
\boxed{\text{This is a successful HTTP response carrying an HTML page.}}
\]

\subsection*{Problem 7: HTTP Status Codes}
Classify each HTTP status code.

\[
\begin{array}{c|c}
\text{Status Code} & \text{Meaning} \\
\hline
200 & ? \\
301 & ? \\
404 & ? \\
500 & ?
\end{array}
\]

\subsection*{Solution}

\[
\begin{array}{c|c}
\text{Status Code} & \text{Meaning} \\
\hline
200 & \text{Success} \\
301 & \text{Redirect} \\
404 & \text{Client error, resource not found} \\
500 & \text{Server error}
\end{array}
\]

\[
\boxed{
\begin{array}{c}
2xx = \text{success} \\
3xx = \text{redirect} \\
4xx = \text{client error} \\
5xx = \text{server error}
\end{array}
}
\]

\subsection*{Problem 8: HTTP/1.1 Head-of-Line Blocking}
Why did HTTP/1.1 pipelining fail in practice?

\subsection*{Solution}
HTTP/1.1 pipelining allowed multiple requests to be sent on one TCP connection.

However, responses had to come back in the same order as the requests.

Suppose the client sends:

\[
R_1, R_2, R_3
\]

If response \(R_1\) is slow, then \(R_2\) and \(R_3\) must wait, even if they are ready.

\[
R_1 \text{ slow} \Rightarrow R_2 \text{ blocked} \Rightarrow R_3 \text{ blocked}
\]

\[
\boxed{\text{One slow response blocks every response behind it.}}
\]

This is called application-layer head-of-line blocking.

\subsection*{Problem 9: HTTP/2 Multiplexing}
How does HTTP/2 improve on HTTP/1.1?

\subsection*{Solution}
HTTP/2 keeps the same HTTP meaning, such as:

\[
\text{methods, headers, status codes}
\]

but changes the wire format.

HTTP/2 uses:

\begin{enumerate}
    \item Binary framing
    \item Multiplexing
    \item Header compression
\end{enumerate}

With multiplexing, many streams can share one TCP connection.

\[
\text{Stream 1, Stream 2, Stream 3} \rightarrow \text{One TCP connection}
\]

Frames from different streams can be interleaved.

\[
\boxed{\text{HTTP/2 solves HTTP/1.1 application-layer head-of-line blocking.}}
\]

\subsection*{Problem 10: HTTP/3 and QUIC}
Why does HTTP/3 use QUIC over UDP instead of TCP?

\subsection*{Solution}
HTTP/2 still runs over TCP. Even though HTTP/2 multiplexes streams, TCP delivers bytes in order.

If one TCP segment is lost, all streams must wait.

\[
\text{Lost TCP segment} \Rightarrow \text{all HTTP/2 streams stall}
\]

HTTP/3 uses QUIC, which runs over UDP.

QUIC provides independent streams, so if one stream has packet loss, the others can continue.

\[
\text{Lost packet in Stream 1} \not\Rightarrow \text{Stream 2 blocked}
\]

\[
\boxed{\text{HTTP/3 removes TCP-layer head-of-line blocking.}}
\]

\section*{Mixed DNS and HTTP Problems}

\subsection*{Problem 11: Loading a Web Page}
A browser loads:

\[
\texttt{www.example.com}
\]

Explain what happens first: DNS or HTTP?

\subsection*{Solution}
DNS happens first.

The browser needs an IP address before it can open a connection to the server.

The order is:

\begin{enumerate}
    \item Browser asks DNS for the IP address of \texttt{www.example.com}
    \item DNS returns an IP address
    \item Browser opens a transport connection
    \item Browser sends an HTTP request
    \item Server returns an HTTP response
\end{enumerate}

\[
\boxed{\text{DNS finds the server. HTTP asks the server for the resource.}}
\]

\subsection*{Problem 12: Browser as Coordinator}
HTTP fetches one resource at a time. Explain why loading one web page can create many HTTP requests.

\subsection*{Solution}
The first HTTP request usually fetches an HTML document.

That HTML document may reference many other resources:

\[
\text{CSS, JavaScript, images, videos, fonts}
\]

The browser parses the HTML and sends additional HTTP requests for those resources.

For example:

\[
1 \text{ HTML file} + 20 \text{ images} + 3 \text{ CSS files} + 5 \text{ JS files}
\]

\[
= 29 \text{ HTTP requests}
\]

\[
\boxed{\text{The browser coordinates many HTTP fetches to build one web page.}}
\]

\section*{Quick Review}

\[
\boxed{
\begin{array}{l}
\text{DNS maps names to IP addresses.} \\
\text{HTTP fetches resources from servers.} \\
\text{HTTP/1.1 is text-based and uses keep-alive.} \\
\text{HTTP/2 adds binary framing and multiplexing.} \\
\text{HTTP/3 uses QUIC over UDP to avoid TCP head-of-line blocking.}
\end{array}
}
\]

\end{document}